% GENERATED FILE. Edit canonical lessons, then run npm run generate:books.
% canonical-source-hash: fnv1a64:bc866efb4cd57890
% canonical-lessons: PT-C283-boato, PT-C283-rumor, PT-C283-desmentir, PT-C283-desmentido, PT-C283-porta-voz

\chapter{Boatos e desmentidos}
\label{ch:pt-boatos-e-desmentidos}
\hlchaptermodality{283}

\begin{chapteropening}
\textbf{By the end of this chapter:} \emph{I can talk about rumours and denials, and name the spokesperson.}

Complete the last lesson of chapter 283: I can talk about rumours and denials, and name the spokesperson.
\end{chapteropening}

\section[o boato]{o boato — the rumour}

\label{lesson:PT-C283-boato}



\begin{quote}
Before the new one: say the Portuguese for the press release, then the Portuguese for the commentator.
\end{quote}

\subsection*{You'll want to know: o boato}
\textbf{o boato} — \textquotedblleft{}the rumour\textquotedblright{}.

\begin{quote}
Corre um boato na escola: o diretor vai sair.
\end{quote}

\begin{quote}
\emph{There is a rumour going round the school: the head is leaving.}
\end{quote}

\emph{Corre um boato} — a rumour is going round.

\begin{cousinweb}
\emph{boato} is a learned borrowing from Latin \emph{boātus} \textquotedblleft{}a bellowing, a loud cry\textquotedblright{} (EtymDB), from \emph{boāre} \textquotedblleft{}to cry aloud\textquotedblright{}.
\end{cousinweb}

\begin{grammarlens}[title={What you've built}]
One more word for telling, explaining and describing at B1.
\end{grammarlens}

\subsection*{Your turn}
\begin{itemize}
\raggedright
  \item \emph{Say it:} \emph{o boato}
  \item \emph{Say it:} \emph{o boato}, once more
  \item \emph{Say it:} a sentence of your own with \emph{o boato}
  \item \emph{Recall:} say the Portuguese for the press release, then the Portuguese for the commentator, then say \emph{o boato} again
\end{itemize}

\begin{tcolorbox}[breakable,colback=teal!4,colframe=teal!35!black,title={Before you move on}]
What does o boato mean? (\textquotedblleft{}The rumour\textquotedblright{}.) Say it once more.
\end{tcolorbox}
\section[o rumor]{o rumor — the murmur, a low noise; a rumour}

\label{lesson:PT-C283-rumor}



\begin{quote}
Before the new one: say the Portuguese for the commentator, then the Portuguese for the rumour.
\end{quote}

\subsection*{You'll want to know: o rumor}
\textbf{o rumor} — \textquotedblleft{}the murmur, a low noise; a rumour\textquotedblright{}.

\begin{quote}
Ouvia-se o rumor do mar.
\end{quote}

\begin{quote}
\emph{You could hear the murmur of the sea.}
\end{quote}

In news it can also mean a rumour, like \emph{o boato}.

\begin{cousinweb}
\emph{rumor} comes from Latin \emph{rūmor} \textquotedblleft{}a noise, common talk, a rumour\textquotedblright{} (EtymDB).
\end{cousinweb}

\begin{grammarlens}[title={What you've built}]
One more word for telling, explaining and describing at B1.
\end{grammarlens}

\subsection*{Your turn}
\begin{itemize}
\raggedright
  \item \emph{Say it:} \emph{o rumor}
  \item \emph{Say it:} \emph{o rumor}, once more
  \item \emph{Say it:} a sentence of your own with \emph{o rumor}
  \item \emph{Recall:} say the Portuguese for the commentator, then the Portuguese for the rumour, then say \emph{o rumor} again
\end{itemize}

\begin{tcolorbox}[breakable,colback=teal!4,colframe=teal!35!black,title={Before you move on}]
What does o rumor mean? (\textquotedblleft{}The murmur, a low noise; a rumour\textquotedblright{}.) Say it once more.
\end{tcolorbox}
\section[desmentir]{desmentir — to deny, to refute}

\label{lesson:PT-C283-desmentir}



\begin{quote}
Before the new one: say the Portuguese for the rumour, then the Portuguese for the murmur.
\end{quote}

\subsection*{You'll want to know: desmentir}
\textbf{desmentir} — \textquotedblleft{}to deny, to refute\textquotedblright{}.

\begin{quote}
O ator desmentiu a notícia.
\end{quote}

\begin{quote}
\emph{The actor denied the news.}
\end{quote}

Present: \emph{desminto, desmentes, desmente}.

\begin{cousinweb}
\emph{desmentir} is built from \emph{des-} \textquotedblleft{}un-\textquotedblright{} and \emph{mentir} \textquotedblleft{}to lie\textquotedblright{} (EtymDB has the word), which comes from Latin \emph{mentīrī} \textquotedblleft{}to lie\textquotedblright{} (Lewis \& Short): to show that a story is false.
\end{cousinweb}

\begin{grammarlens}[title={What you've built}]
One more word for telling, explaining and describing at B1.
\end{grammarlens}

\subsection*{Your turn}
\begin{itemize}
\raggedright
  \item \emph{Say it:} \emph{desmentir}
  \item \emph{Say it:} \emph{desmentir}, once more
  \item \emph{Say it:} a sentence of your own with \emph{desmentir}
  \item \emph{Recall:} say the Portuguese for the rumour, then the Portuguese for the murmur, then say \emph{desmentir} again
\end{itemize}

\begin{tcolorbox}[breakable,colback=teal!4,colframe=teal!35!black,title={Before you move on}]
What does desmentir mean? (\textquotedblleft{}To deny, to refute\textquotedblright{}.) Say it once more.
\end{tcolorbox}
\section[o desmentido]{o desmentido — the denial}

\label{lesson:PT-C283-desmentido}



\begin{quote}
Before the new one: say the Portuguese for the murmur, then the Portuguese for to deny.
\end{quote}

\subsection*{You'll want to know: o desmentido}
\textbf{o desmentido} — \textquotedblleft{}the denial\textquotedblright{}.

\begin{quote}
O clube respondeu com um desmentido no mesmo dia.
\end{quote}

\begin{quote}
\emph{The club replied with a denial the same day.}
\end{quote}

The past participle of \emph{desmentir} used as a noun.

\begin{cousinweb}
\emph{desmentido} is the past participle of \emph{desmentir} \textquotedblleft{}to deny\textquotedblright{} used as a noun; the verb is built on \emph{des-} and a verb that goes back to Latin \emph{mentīrī} \textquotedblleft{}to lie\textquotedblright{} (Lewis \& Short).
\end{cousinweb}

\begin{grammarlens}[title={What you've built}]
One more word for telling, explaining and describing at B1.
\end{grammarlens}

\subsection*{Your turn}
\begin{itemize}
\raggedright
  \item \emph{Say it:} \emph{o desmentido}
  \item \emph{Say it:} \emph{o desmentido}, once more
  \item \emph{Say it:} a sentence of your own with \emph{o desmentido}
  \item \emph{Recall:} say the Portuguese for the murmur, then the Portuguese for to deny, then say \emph{o desmentido} again
\end{itemize}

\begin{tcolorbox}[breakable,colback=teal!4,colframe=teal!35!black,title={Before you move on}]
What does o desmentido mean? (\textquotedblleft{}The denial\textquotedblright{}.) Say it once more.
\end{tcolorbox}
\section[o porta-voz]{o porta-voz — the spokesperson}

\label{lesson:PT-C283-porta-voz}



\begin{quote}
Before the new one: say the Portuguese for to deny, then the Portuguese for the denial.
\end{quote}

\subsection*{You'll want to know: o porta-voz}
\textbf{o porta-voz} — \textquotedblleft{}the spokesperson\textquotedblright{}.

\begin{quote}
O porta-voz do governo falou à imprensa.
\end{quote}

\begin{quote}
\emph{The government spokesperson spoke to the press.}
\end{quote}

The same form for a woman: \emph{a porta-voz}.

\begin{cousinweb}
\emph{porta-voz} joins \emph{porta} \textquotedblleft{}carries\textquotedblright{}, from Latin \emph{portāre} \textquotedblleft{}to carry\textquotedblright{}, and \emph{voz} \textquotedblleft{}voice\textquotedblright{}, from Latin \emph{vōx} (Lewis \& Short): the one who carries the voice.
\end{cousinweb}

\begin{grammarlens}[title={What you've built}]
One more word for telling, explaining and describing at B1.
\end{grammarlens}

\subsection*{Your turn}
\begin{itemize}
\raggedright
  \item \emph{Say it:} \emph{o porta-voz}
  \item \emph{Say it:} \emph{o porta-voz}, once more
  \item \emph{Say it:} a sentence of your own with \emph{o porta-voz}
  \item \emph{Recall:} say the Portuguese for to deny, then the Portuguese for the denial, then say \emph{o porta-voz} again
\end{itemize}

\begin{tcolorbox}[breakable,colback=teal!4,colframe=teal!35!black,title={Before you move on}]
What does o porta-voz mean? (\textquotedblleft{}The spokesperson\textquotedblright{}.) Say it once more.
\end{tcolorbox}
